\section{Demonstration}
\label{sec:demo}

\begin{figure*}[t]
  \centering
  \includegraphics[width=0.88\textwidth]{alembic_interface.png}
  \caption{The Alembic web dashboard. The top bar accepts a source repository URL
  and controls execution; the stage rail beneath it visualizes progress through
  the explorer $\rightarrow$ environment $\rightarrow$ coder $\rightarrow$
  validator $\rightarrow$ wrapper stages. The left panel shows the exploration
  report, the generated output files, how to run them (the recorded
  \texttt{setup.sh}), and per-tool invocation examples; the right panel shows the
  generated MCP server together with its validation badges. A collapsible
  activity feed at the bottom streams individual steps in real time, and
  completed tools can be invoked interactively with user-supplied arguments.}
  \label{fig:dashboard}
\end{figure*}

Alembic targets engineers building agent-based scientific assistants who need a
growing tool library without hand-writing a wrapper per repository, and
maintainers wanting a low-effort path to agent-accessible software. To make the
pipeline observable, we provide a lightweight web dashboard
(Figure~\ref{fig:dashboard}) that streams execution in real time and exposes the
synthesized MCP server. As the five stages unfold live, a stage rail tracks the
active agent's success or failure and a feed logs each step. Rather than raw
output, the dashboard shows pass/fail badges for static checks and test runs,
and tool invocations annotated with inputs, outputs, and errors; any generated
tool can be invoked with custom arguments and its result seen immediately. The
screencast (Section~\ref{sec:availability}) walks a full run from a bare GitHub
URL to a served MCP, and shows the \texttt{--resume <stage>} flag and batch mode
(\texttt{run\_benchmark.py}) over many repositories.